% IEEE Access draft — Adaptive Parallel Elastics for Energy-Efficient Quadruped Load Carriage
% Port to ieeeaccess.cls for submission; uses IEEEtran here for portability.
\documentclass[journal]{IEEEtran}
\usepackage{graphicx,amsmath,booktabs,url,hyperref}
\begin{document}

\title{Gearing Decides the Sign: An Adaptive, Self-Tuning Parallel Knee Preload for
Energy-Efficient Quadruped Load Carriage}

\author{I.~Osokin \emph{et al.}% TODO author list/affiliations
\thanks{Preprint draft, \today. Simulation study; absolute electrical figures are reported as bands
pending measured motor constants.}}

\maketitle

\begin{abstract}
Adding a passive parallel elastic element beside a motor offloads joint torque and, because resistive
(ohmic) motor loss scales with the \emph{square} of torque, should cut electrical energy. Whether it does,
however, is governed by the actuator's gear ratio, which sets how large a share of the electrical budget the
ohmic term occupies. We study this trade-off end-to-end in simulation (MuJoCo Playground / MJX, brax PPO) on
two commercial Unitree platforms. On the high-gear G1 humanoid (ohmic $\approx$4\% of the budget) an in-loop
parallel hip spring makes walking \emph{worse} (+7\%), reversing an optimistic post-hoc estimate. On the
low-gear Go1 quadruped (ohmic $\approx$54\%) a \emph{constant} parallel knee preload cuts the cost of transport
(CoT) by 14--27\% across three of four training conditions, with the benefit growing with payload. We further
introduce an \emph{adaptive, per-leg, self-tuning} preload: a slow outer loop that senses each knee's own motor
torque and ramps a passive preload to offload it, using no load sensor and no payload observation; it carries
the energy win into blind load-carrying (0--6\,kg). We report the boundary where the mechanism reverses, the
constant-preload (vs.\ linear-spring) element insight for the offset-dominated knee, and the high-load
stability cost honestly. The result is a deployable, sensor-free efficiency mechanism for low-gear load-carrying
quadrupeds, and a clear statement of when parallel elasticity helps and when it hurts.
\end{abstract}

\begin{IEEEkeywords}
parallel elastic actuator, cost of transport, legged robots, quadruped, reinforcement learning, energy efficiency.
\end{IEEEkeywords}

\section{Introduction}
Electric legged robots spend much of their energy as heat in the motor windings. For a geared joint the
copper loss is $P_\Omega = (\tau/K_t)^2 R$, quadratic in joint torque $\tau$, so any mechanism that offloads
torque from the motor cuts heat super-linearly. A \emph{parallel} elastic element---a spring beside the
motor---does exactly this for the quasi-static support torque a leg holds during stance, and additionally stores
and returns energy the motor would otherwise dissipate (most commercial drivers do not regenerate).

The catch is twofold. First, a parallel spring is \emph{always engaged}: it also resists the motor in phases
where its torque is unwanted, so the policy must learn to exploit it. Second, and decisively, the ohmic term is
only worth attacking if it is a large share of the budget---and that share is set by the \textbf{gear ratio}.
High reduction makes the motor-side torque small and the ohmic loss negligible; low reduction leaves large
torque at the motor, where offloading pays.

\textbf{Contributions.}
(1) We show \emph{gearing decides the sign}: the same parallel-spring idea is harmful on a high-gear humanoid
and beneficial on a low-gear quadruped, quantified in matched in-loop experiments.
(2) For the low-gear Go1 quadruped we measure a \textbf{14--27\% cost-of-transport reduction} from a
\emph{constant} parallel knee preload, and show the element must be a constant preload, not a linear
spring---the knee work-loop is offset-dominated.
(3) We introduce an \textbf{adaptive, per-leg, self-tuning preload} that requires no load sensor and no payload
observation, and carries the energy win into blind load-carrying with the benefit \emph{growing with load}.
(4) We report the limits honestly: a seed-dependent magnitude, a high-load stability cost, and a sim-to-real
realism bound on the payload range.

\section{Related Work}
\label{sec:related}
% TODO: enrich from docs/literature_review.md (autonomous lit-review in progress) — SEA/PEA, VSA, clutched.
Parallel-elastic torque offloading to cut motor energy is foundational~\cite{plooij2012}. Torque-squared
(electrical) accounting for a parallel spring, validated on hardware, has been demonstrated on a
quadruped~\cite{bjelonic2023}. Reinforcement learning with elastic actuators for legged locomotion has been
explored on quadrupeds~\cite{bjelonic2022}. Tunable parallel springs optimized with a torque-squared metric
were studied analytically on a leg test-stand~\cite{belov2024}, without RL, gait, or a load-adaptive loop.
Our work differs by (i) co-adapting the gait to the element in-loop on commercial platforms, (ii) mapping the
gear-ratio boundary where the benefit reverses, and (iii) the sensor-free, per-leg, load-adaptive preload.
A fuller survey of series vs.\ parallel elasticity, variable-stiffness actuators, clutched/quasi-passive
elements, and biological running springs is given in the supplementary literature review.

\section{Method}
\subsection{Platforms, simulator, and policy}
We use MuJoCo Playground's Unitree Go1 (12\,kg, 6.33:1 reduction) and G1 (high reduction) joystick environments,
simulated with MJX and trained with brax PPO (8192 parallel environments, 200\,M steps). The policy is a
feed-forward actor--critic tracking a commanded body velocity at 50\,Hz, under standard domain randomization.

\subsection{Parallel spring model and honest energy accounting}
The spring torque $\tau_{\mathrm{spring}}(\theta)$ is injected at the target joint through the external
generalized force \texttt{qfrc\_applied}, \emph{not} through the actuators, so the motor torque
\texttt{qfrc\_actuator} excludes the passive contribution and the electrical accounting stays honest. Per
actuated degree of freedom the electrical power is
\begin{equation}
P_{\mathrm{elec}} = \max\!\big(\tau\,\omega + (\tau/K_t)^2 R,\; 0\big),
\end{equation}
i.e.\ mechanical plus ohmic with \emph{no regeneration} (the $\max(\cdot,0)$ clamp), applied identically in the
training reward and the evaluation metric. The headline metric is cost of transport, $\mathrm{CoT} =
P_{\mathrm{elec}}/v$ (electrical watts per metre per second), which normalizes for the speed differences the
spring induces.

\subsection{The element: a constant preload, not a linear spring}
Fitting a linear spring to the Go1 knee work-loop degenerates to $k\approx 0$: the loop is dominated by a
near-constant support offset through stance, not a restoring slope (a $\tau^2$ fit is even mildly
anti-restoring). The buildable optimum is therefore a \emph{constant preload}, realizable as a heavily
pre-wound low-rate coil whose torque is nearly constant over the knee's small operating range.

\subsection{Adaptive, per-leg, self-tuning preload}
A slow outer loop wraps the fast policy. For each leg $i$ it maintains a 15\,s exponential moving average
$\bar\tau_i$ of that knee's \emph{own} motor torque and ramps the passive preload toward full compensation,
\begin{equation}
\dot\tau_{0,i} = \mathrm{clip}\!\big(k_p\,\bar\tau_i,\, -2,\, +2\big)\ \mathrm{N\,m/s},\quad k_p \approx 0.2,
\end{equation}
with no load sensor and no payload observation. The four legs are independent reflex loops, so front/rear and
left/right asymmetry is handled automatically. Because training episodes are shorter than the 15\,s window, we
\emph{train robust} (preload domain-randomized per episode) and \emph{adapt at deploy} (run the controller in
long rollouts so $\tau_0$ converges); the converged $\tau_0$ scales with load and is larger on the front legs.

\section{Experiments and Results}
\subsection{Gearing decides the sign}
On the high-gear G1, where ohmic loss is $\approx$4\% of the motor budget, an in-loop parallel hip spring makes
walking \textbf{+7\% worse} and less stable, reversing an optimistic post-hoc $-3.8\%$ (the always-on spring
fights the gait with no benefit to recover). On the low-gear Go1, ohmic is $\approx$54\% of the budget and the
knee preload pays (below). This sign flip with gear ratio is the central organizing result.

\subsection{Go1 load-carrying: cost of transport vs.\ load}
Table~\ref{tab:cot} reports the CoT reduction of the adaptive preload against a matched no-spring baseline on
flat ground, for three seeds and a payload curriculum; Fig.~\ref{fig:cot} plots CoT and the reduction vs.\ load.

\begin{table}[t]\centering
\caption{Cost-of-transport reduction (\%), adaptive knee preload vs.\ matched no-spring baseline, flat ground.}
\label{tab:cot}
\begin{tabular}{lrrr}
\toprule
Condition & @0\,kg & @2.5\,kg & @5\,kg \\
\midrule
Seed 1 & $-16.6$ & $-19.5$ & $-22.8$ \\
Seed 2 (weak outlier) & $-3.4$ & $-8.3$ & $-6.5$ \\
Seed 3 & $-13.9$ & $-20.1$ & $-26.7$ \\
Payload curriculum & $-16.8$ & $-20.4$ & $-22.0$ \\
\bottomrule
\end{tabular}
\end{table}

\begin{figure}[t]\centering
\includegraphics[width=\columnwidth]{../outputs/figures/cot_vs_load.png}
\caption{Cost of transport (left; solid = baseline, dashed = adaptive) and its reduction (right) vs.\ payload.
The benefit grows with load in three of four conditions; seed 2 is a weak outlier.}
\label{fig:cot}
\end{figure}

The adaptive preload cuts CoT by \textbf{14--27\% in three of four conditions}, and the reduction \emph{grows
with payload}---heavier loads raise the support torque the preload offloads. One seed (seed~2) is a weak
$-3$ to $-8\%$ outlier: with the offload available, that policy spends it on walking faster rather than on
energy, which the speed-normalized CoT only partly compensates.

\section{Discussion and Limitations}
We foreground the threats to validity rather than bury them.
\textbf{(1) Stability cost.} At high load ($\geq$7.5\,kg) the adaptive policy loses some survival (e.g.\ 1070/1500
at 10\,kg) where the matched baseline holds 1500/1500; the energy win is a low-to-mid-load result.
\textbf{(2) Realism bound.} The real Go1 carries $\sim$5--10\,kg; the simulator enforces peak but not
continuous/thermal torque, structure, or balance, so high-load (15--30\,kg) ``walking'' is sim-only. The
physically meaningful range is 0--6\,kg; the 15--30\,kg regime would belong to a B2-class robot.
\textbf{(3) No-regeneration dependence.} The win lives partly in dissipated braking energy; with regeneration it
shrinks ($\sim$24\% sensitivity). No-regen is justified by back-EMF physics for these drivers but is not a
measured spec.
\textbf{(4) Absolute energy is a band.} We lack measured motor $K_t,R$; relative CoT comparisons are robust, but
absolute watts are uncertain.
\textbf{(5) Energy-naive baseline and rough terrain.} The baseline does not itself optimize energy hard, and on
rough (2.5\,cm) terrain both arms are unstable ($\sim$40\% survival), so the win there is inconclusive.

\section{Conclusion}
Whether a parallel elastic element saves electrical energy in a legged robot is decided by the gear ratio: it is
harmful on a high-gear humanoid and beneficial on a low-gear quadruped. For the low-gear Go1 we demonstrate a
14--27\% cost-of-transport reduction from a constant knee preload, and a sensor-free, per-leg, self-tuning
adaptive preload that carries the win into blind load-carrying with the benefit growing with load. The mechanism
is simple and deployable; its boundary, its element form, and its high-load stability cost are characterized.
Future work runs the same idea on a \emph{running} quadruped, where impact braking energy is largest and a
one-sided stiffness spring may add a recovery channel a constant preload cannot.

\begin{thebibliography}{1}
\bibitem{plooij2012} M.~Plooij and M.~Wisse, ``A novel spring mechanism to reduce energy consumption of robotic
arms,'' \emph{IEEE/RSJ IROS}, 2012.
\bibitem{bjelonic2023} F.~Bjelonic \emph{et al.}, ``Learning-based design and control for quadrupedal robots with
parallel-elastic actuators,'' \emph{IEEE RA-L}, 2023.
\bibitem{bjelonic2022} F.~Bjelonic \emph{et al.}, ``Reinforcement learning for elastic quadrupedal locomotion,''
2022.
\bibitem{belov2024} Belov, Osokin \emph{et al.}, ``Torque-squared-optimal parallel elastic elements for a robotic
leg,'' Skoltech, 2024.
% TODO: add references from docs/literature_review.md
\end{thebibliography}

\end{document}
